\section{Discussion}

\subsection{Advantages and limitations by method family}

\begin{table}[H]
\centering
\caption{Advantages and limitations by method family.}
\small
\begin{tabular}{@{}p{3.4cm}p{5.5cm}p{5.5cm}@{}}
\toprule
Approach & Advantages & Limitations \\
\midrule
TF-IDF + classical head & Strong on lexical topics; fast; interpretable; HDF5 persistence & Weaker on compositional polarity; stop-word removal can hurt intent/sentiment \\
Frozen BGE-small + LogReg & Contrastive sentence geometry; strong on RT/Banking77; calibrated LogReg & Requires Transformer runtime; not task-fine-tuned; fixed encoder \\
TextCNN (bonus) & End-to-end local patterns; probability outputs for Brier/ECE & Needs more labelled data; lags on 77-way Banking77 in our runs \\
\bottomrule
\end{tabular}
\end{table}

\subsection{Learning outcomes relative to course concepts}

The three-corpus design confirms that ``best model'' is not universal. Conditional independence (Naive Bayes) remains a useful sparse baseline on AG News. Margin and logistic linear heads excel when the feature space already encodes discriminative patterns---whether sparse $n$-grams or frozen BGE vectors. Frozen transfer via a \emph{contrastive} sentence encoder differs from using an MLM encoder as a bag feature: BGE is trained for sentence-level similarity, which aligns better with short-query intent and sentiment. TextCNN illustrates that a small convolutional network can learn local patterns from scratch but does not automatically dominate classical or frozen-embedding pipelines at this labelled-data scale. Finally, calibration and negation splits show that Macro-F1 alone is an incomplete academic evaluation.
